\hypertarget{python-3.5-native-asyncawait-coroutines-and-matrix-operations}{%
\chapter{Python 3.5: Native Async/Await Coroutines and Matrix
Operations}\label{python-3.5-native-asyncawait-coroutines-and-matrix-operations}}

\hypertarget{pep-492-native-coroutines-asyncawait-internals}{%
\subsection{\texorpdfstring{11.1 PEP 492: Native Coroutines
(\texttt{async/await}
Internals)}{11.1 PEP 492: Native Coroutines (async/await Internals)}}\label{pep-492-native-coroutines-asyncawait-internals}}

Python 3.5 introduced native coroutines via PEP 492, establishing a
clear separation between standard generators and asynchronous tasks to
prevent logical developer errors.

\hypertarget{the-generator-coroutine-limitation}{%
\subsubsection{1. The Generator-Coroutine
Limitation}\label{the-generator-coroutine-limitation}}

In Python 3.4, coroutines were decorated generators
(\passthrough{\lstinline!@asyncio.coroutine!}). Under the hood, these
coroutines were instances of \passthrough{\lstinline!PyGenObject!}.
Because they shared the same C type as standard generators, they exposed
the standard iterator protocol. Developers could mistakenly iterate over
a coroutine using a \passthrough{\lstinline!for!} loop or call
\passthrough{\lstinline!next()!} on it directly, causing unexpected
runtime behavior. Additionally, this lack of syntactic separation meant
that: * The interpreter could not perform static analysis to detect
missing
\passthrough{\lstinline!await!}/\passthrough{\lstinline!yield from!}
calls on coroutine objects. * It was possible to call
\passthrough{\lstinline!next(coro())!} directly, stepping the generator
and bypassing the event loop's state machine. * The
\passthrough{\lstinline!yield from!} syntax was overloaded, being used
for both yielding from standard generators and awaiting asynchronous
tasks, creating a high degree of cognitive load and making code
refactoring prone to errors.

\hypertarget{cpython-representation-pycoroobject}{%
\subsubsection{\texorpdfstring{2. CPython Representation:
\texttt{PyCoroObject}}{2. CPython Representation: PyCoroObject}}\label{cpython-representation-pycoroobject}}

Native coroutines declared using \passthrough{\lstinline!async def!}
compile directly to a specialized C structure,
\passthrough{\lstinline!PyCoroObject!}
(\passthrough{\lstinline!Include/genobject.h!}):

\begin{lstlisting}[language=C]
typedef struct {
    PyObject_HEAD
    struct _frame *cr_frame;     /* Frame object executing coroutine code */
    PyObject *cr_code;          /* Code object compiled from async def */
    PyObject *cr_name;          /* Coroutine name */
    PyObject *cr_qualname;      /* Qualified name */
    PyObject *cr_origin;        /* Awaited task or future origin trace */
    PyObject *cr_weakreflist;   /* Weak reference list pointer */
    char cr_running;            /* Flag indicating if coroutine is active */
} PyCoroObject;
\end{lstlisting}

Let's break down the role of each field in the CPython runtime: *
\textbf{\passthrough{\lstinline!PyObject\_HEAD!}}: The standard object
header containing the reference counter
(\passthrough{\lstinline!ob\_refcnt!}) and the pointer to the type
object (\passthrough{\lstinline!ob\_type!}). In this case,
\passthrough{\lstinline!ob\_type!} points to the
\passthrough{\lstinline!PyCoro\_Type!} struct. *
\textbf{\passthrough{\lstinline!cr\_frame!}}: A pointer to the execution
frame (\passthrough{\lstinline!PyFrameObject!}). It holds the local
execution state, the bytecode pointer, the local evaluation stack, and
local variable cells. * \textbf{\passthrough{\lstinline!cr\_code!}}:
Points to the \passthrough{\lstinline!PyCodeObject!} representing the
compiled bytecode. * \textbf{\passthrough{\lstinline!cr\_name!}}: The
coroutine's name (represented by a
\passthrough{\lstinline!PyUnicodeObject!}). *
\textbf{\passthrough{\lstinline!cr\_qualname!}}: The qualified name
(e.g.~\passthrough{\lstinline!ClassName.method\_name!}). *
\textbf{\passthrough{\lstinline!cr\_origin!}}: If debugging is enabled
via
\passthrough{\lstinline!sys.set\_coroutine\_origin\_tracking\_depth()!},
this pointer stores a traceback tuple showing where the coroutine was
instantiated, enabling deep debugging of orphaned or un-awaited
coroutines. * \textbf{\passthrough{\lstinline!cr\_weakreflist!}}: Head
of a doubly-linked list of weak reference objects pointing to this
coroutine. * \textbf{\passthrough{\lstinline!cr\_running!}}: A
single-byte boolean flag. If \passthrough{\lstinline!cr\_running!} is
true, it indicates the coroutine frame is currently executing. Any
attempt to resume the coroutine from another context will raise a
\passthrough{\lstinline!RuntimeError: coroutine already running!},
protecting the runtime from re-entrant frame execution.

Because \passthrough{\lstinline!PyCoroObject!} is a distinct C type, it
does not implement sequence or iterator slots
(\passthrough{\lstinline!tp\_iter!} is NULL), preventing standard
iteration errors.

\hypertarget{the-am_await-slot-protocol}{%
\subsubsection{\texorpdfstring{3. The \texttt{am\_await} Slot
Protocol}{3. The am\_await Slot Protocol}}\label{the-am_await-slot-protocol}}

Instead of iteration slots, \passthrough{\lstinline!PyCoroObject!}
implements the \passthrough{\lstinline!tp\_as\_async!} protocol table,
specifically the \passthrough{\lstinline!am\_await!} slot:

\begin{lstlisting}[language=C]
typedef struct {
    unaryfunc am_await;         /* __await__ implementation pointer */
    unaryfunc am_aiter;         /* __aiter__ implementation pointer */
    unaryfunc am_anext;         /* __anext__ implementation pointer */
} PyAsyncMethods;
\end{lstlisting}

When an object is awaited via \passthrough{\lstinline!await expr!}, the
CPython interpreter executes the following internal evaluation loop
logic: 1. It checks if the object's type has
\passthrough{\lstinline!tp\_as\_async!} populated and if
\passthrough{\lstinline!tp\_as\_async->am\_await!} is non-NULL. 2. If
yes, it calls \passthrough{\lstinline!am\_await(expr)!}. This slot
function MUST return an iterator (an object that implements
\passthrough{\lstinline!tp\_iternext!}). 3. For native coroutines
(\passthrough{\lstinline!PyCoroObject!}), their type has a custom
\passthrough{\lstinline!am\_await!} slot that wraps the coroutine in a
\passthrough{\lstinline!PyCoroWrapper!} (or returns a wrapper that
exposes standard iterator operations for the event loop to drive). 4.
For user-defined classes, a custom
\passthrough{\lstinline!\_\_await\_\_!} method is defined. At the C
level, this maps to \passthrough{\lstinline!am\_await!} resolving the
python-level \passthrough{\lstinline!\_\_await\_\_!} method. It must
return an iterator.

\begin{lstlisting}[language=Python]
# User-level custom awaitable implementation
class DatabaseConnection:
    def __await__(self):
        # We yield control back to the event loop if the socket is not ready
        while not self.is_connected():
            yield None
        return self._connection
\end{lstlisting}

\hypertarget{bytecode-compilation-and-tracing-get_awaitable}{%
\subsubsection{\texorpdfstring{4. Bytecode Compilation and Tracing:
\texttt{GET\_AWAITABLE}}{4. Bytecode Compilation and Tracing: GET\_AWAITABLE}}\label{bytecode-compilation-and-tracing-get_awaitable}}

When the compiler parses an \passthrough{\lstinline!await expression!}
statement, it emits a specialized
\textbf{\passthrough{\lstinline!GET\_AWAITABLE!}} bytecode instruction:
1. \passthrough{\lstinline!GET\_AWAITABLE!} pops the expression result
off the evaluation stack. 2. The VM checks if the object implements the
\passthrough{\lstinline!am\_await!} slot (either natively or via class
overrides). 3. If valid, the VM calls the slot to retrieve the awaitable
iterator and pushes it onto the stack. If not valid, it raises a
\passthrough{\lstinline!TypeError!}. 4. This is followed by a delegation
loop (similar to \passthrough{\lstinline!yield from!}) that yields
control back to the event loop if the awaitable is pending.

Let's trace the compiled bytecodes of an asynchronous execution path.
Consider the following code:

\begin{lstlisting}[language=Python]
async def get_val():
    return 42

async def main():
    val = await get_val()
\end{lstlisting}

The CPython 3.5 compiler generates the following disassemblies:

\hypertarget{disassembly-of-get_val}{%
\paragraph{\texorpdfstring{Disassembly of
\texttt{get\_val}:}{Disassembly of get\_val:}}\label{disassembly-of-get_val}}

\begin{lstlisting}
  2           0 LOAD_CONST               1 (42)
              3 RETURN_VALUE
\end{lstlisting}

\hypertarget{disassembly-of-main}{%
\paragraph{\texorpdfstring{Disassembly of
\texttt{main}:}{Disassembly of main:}}\label{disassembly-of-main}}

\begin{lstlisting}
  5           0 LOAD_GLOBAL              0 (get_val)
              3 CALL_FUNCTION            0
              6 GET_AWAITABLE            0
              9 LOAD_CONST               0 (None)
             12 YIELD_FROM               0
             15 STORE_FAST               0 (val)
             18 LOAD_CONST               0 (None)
             21 RETURN_VALUE
\end{lstlisting}

\hypertarget{step-by-step-vm-execution-trace-of-main}{%
\paragraph{\texorpdfstring{Step-by-Step VM Execution Trace of
\texttt{main()}:}{Step-by-Step VM Execution Trace of main():}}\label{step-by-step-vm-execution-trace-of-main}}

\begin{enumerate}
\def\labelenumi{\arabic{enumi}.}
\tightlist
\item
  \textbf{\passthrough{\lstinline!LOAD\_GLOBAL 0!}}: Resolves the
  identifier \passthrough{\lstinline!get\_val!} from the global
  namespace dictionary and pushes the function object onto the
  evaluation stack.
\item
  \textbf{\passthrough{\lstinline!CALL\_FUNCTION 0!}}: Invokes
  \passthrough{\lstinline!get\_val()!}. Since it was defined with
  \passthrough{\lstinline!async def!}, the VM immediately creates a new
  \passthrough{\lstinline!PyCoroObject!} and pushes it onto the
  evaluation stack. Note that the body of
  \passthrough{\lstinline!get\_val!} does not execute yet.
\item
  \textbf{\passthrough{\lstinline!GET\_AWAITABLE!}}:

  \begin{itemize}
  \tightlist
  \item
    Pops the \passthrough{\lstinline!PyCoroObject!} off the stack.
  \item
    Accesses
    \passthrough{\lstinline!ob\_type->tp\_as\_async->am\_await!}.
  \item
    Executes the slot, which wraps the native coroutine or verifies it
    is ready to be driven.
  \item
    Pushes the resulting awaitable iterator back onto the stack.
  \end{itemize}
\item
  \textbf{\passthrough{\lstinline!LOAD\_CONST 0!}}: Pushes
  \passthrough{\lstinline!None!} onto the stack as the priming value.
\item
  \textbf{\passthrough{\lstinline!YIELD\_FROM!}}:

  \begin{itemize}
  \tightlist
  \item
    Performs a loop that mimics \passthrough{\lstinline!yield from!}
    behavior.
  \item
    It pops the value (\passthrough{\lstinline!None!}) and sends it to
    the awaitable iterator by invoking its
    \passthrough{\lstinline!tp\_iternext!} (or equivalent
    \passthrough{\lstinline!send!} slot).
  \item
    The sub-frame \passthrough{\lstinline!get\_val!} runs and reaches
    \passthrough{\lstinline!RETURN\_VALUE!}. The runtime raises a
    \passthrough{\lstinline!StopIteration!} containing the return value
    \passthrough{\lstinline!42!}.
  \item
    \passthrough{\lstinline!YIELD\_FROM!} catches
    \passthrough{\lstinline!StopIteration!}, extracts
    \passthrough{\lstinline!42!} from the exception's
    \passthrough{\lstinline!value!} attribute, and pushes it onto the
    evaluation stack.
  \end{itemize}
\item
  \textbf{\passthrough{\lstinline!STORE\_FAST 0!}}: Pops the value
  \passthrough{\lstinline!42!} off the stack and stores it in local
  variable \passthrough{\lstinline!val!}.
\end{enumerate}

\begin{lstlisting}
Evaluation Stack Transitions during `await`:

Step 1: [ get_val (fn) ]  <-- LOAD_GLOBAL
Step 2: [ PyCoroObject ]  <-- CALL_FUNCTION
Step 3: [ CoroWrapper ]   <-- GET_AWAITABLE
Step 4: [ CoroWrapper ]   <-- LOAD_CONST (None)
        [    None     ]
Step 5: [     42      ]   <-- YIELD_FROM (catches StopIteration(42))
Step 6: [             ]   <-- STORE_FAST (val = 42)
\end{lstlisting}

\begin{center}\rule{0.5\linewidth}{0.5pt}\end{center}

\hypertarget{matrix-multiplication-operator-pep-465}{%
\subsection{\texorpdfstring{11.2 Matrix Multiplication Operator
(\texttt{@} / PEP
465)}{11.2 Matrix Multiplication Operator (@ / PEP 465)}}\label{matrix-multiplication-operator-pep-465}}

Python 3.5 introduced the binary operator \passthrough{\lstinline!@!}
(and its in-place version \passthrough{\lstinline!@=!}) to support clean
mathematical syntax for matrix multiplication.

\hypertarget{low-level-c-slots-mapping}{%
\subsubsection{1. Low-Level C Slots
Mapping}\label{low-level-c-slots-mapping}}

CPython maps the matrix multiplication operator to two new function
pointer slots inside the \passthrough{\lstinline!PyNumberMethods!}
struct (\passthrough{\lstinline!Include/object.h!}):

\begin{lstlisting}[language=C]
typedef struct {
    /* ... */
    binaryfunc nb_matrix_multiply;          /* Corresponds to __matmul__ */
    binaryfunc nb_inplace_matrix_multiply;  /* Corresponds to __imatmul__ */
} PyNumberMethods;
\end{lstlisting}

When evaluating \passthrough{\lstinline!A @ B!}, CPython's binary
operation dispatch system invokes
\passthrough{\lstinline!PyNumber\_MatrixMultiply(A, B)!}: 1.
\textbf{Left-to-Right Dispatch}: If the type of
\passthrough{\lstinline!A!} defines
\passthrough{\lstinline!tp\_as\_number->nb\_matrix\_multiply!}, it
executes the slot function. 2. \textbf{Right-to-Left Fallback}: If
\passthrough{\lstinline!A!}'s slot is NULL, or it returns
\passthrough{\lstinline!Py\_NotImplemented!}, CPython checks if the type
of \passthrough{\lstinline!B!} defines
\passthrough{\lstinline!nb\_matrix\_multiply!}. 3. \textbf{Subclass
Precedence}: If \passthrough{\lstinline!B!} is a subclass of
\passthrough{\lstinline!A!} and overrides
\passthrough{\lstinline!nb\_matrix\_multiply!},
\passthrough{\lstinline!B!}'s slot is checked and called \emph{before}
\passthrough{\lstinline!A!}'s slot, allowing specialized subclass
multiplication overrides. 4. \textbf{TypeError}: If neither slot returns
a valid object (or both return
\passthrough{\lstinline!Py\_NotImplemented!}), the VM raises a
\passthrough{\lstinline!TypeError: unsupported operand type(s) for @!}.

\hypertarget{python-level-interface-and-numeric-implementations}{%
\subsubsection{2. Python-Level Interface and Numeric
Implementations}\label{python-level-interface-and-numeric-implementations}}

Developers can customize this operator behavior by implementing the
Python magic methods \passthrough{\lstinline!\_\_matmul\_\_!},
\passthrough{\lstinline!\_\_rmatmul\_\_!}, and
\passthrough{\lstinline!\_\_imatmul\_\_!}:

\begin{lstlisting}[language=Python]
class CustomMatrix:
    def __init__(self, grid):
        self.grid = grid

    def __matmul__(self, other):
        if not isinstance(other, CustomMatrix):
            return NotImplemented
        # Compute the dot product matrix
        rows_A, cols_A = len(self.grid), len(self.grid[0])
        rows_B, cols_B = len(other.grid), len(other.grid[0])
        assert cols_A == rows_B, "Dimension mismatch!"
        
        result = [[0] * cols_B for _ in range(rows_A)]
        for i in range(rows_A):
            for j in range(cols_B):
                result[i][j] = sum(self.grid[i][k] * other.grid[k][j] for k in range(cols_A))
        return CustomMatrix(result)
\end{lstlisting}

\hypertarget{high-performance-implementations}{%
\subsubsection{3. High-Performance
Implementations}\label{high-performance-implementations}}

This dedicated operator allowed numerical libraries (like NumPy) to
bypass verbose nested function call chains. Consider this linear algebra
comparison:

\begin{lstlisting}[language=Python]
# Pre-Python 3.5 (Verbose and nesting-heavy)
result = A.dot(B).dot(C) + D.dot(E)

# Python 3.5+ (Clean, mathematical representation)
result = (A @ B @ C) + (D @ E)
\end{lstlisting}

NumPy defines C-level slots mapping for
\passthrough{\lstinline!nb\_matrix\_multiply!} on its
\passthrough{\lstinline!ndarray!} types, executing optimized BLAS or
LAPACK matrix routines underneath. By pushing the matrix multiplication
to these compiled C/Fortran libraries, they release the Global
Interpreter Lock (GIL) during compilation of large arrays, permitting
true multi-threaded parallel computation across CPU cores.

\begin{center}\rule{0.5\linewidth}{0.5pt}\end{center}

\hypertarget{extended-unpacking-generalizations-pep-448}{%
\subsection{11.3 Extended Unpacking Generalizations (PEP
448)}\label{extended-unpacking-generalizations-pep-448}}

PEP 448 expanded the capabilities of the unpacking operators
\passthrough{\lstinline!*!} (iterable unpacking) and
\passthrough{\lstinline!**!} (dictionary unpacking), allowing multiple
unpacks inside collection literals and function calls.

\hypertarget{bytecode-compilation-and-unpacking-instructions}{%
\subsubsection{1. Bytecode Compilation and Unpacking
Instructions}\label{bytecode-compilation-and-unpacking-instructions}}

Prior to Python 3.5, unpacking was limited to a single occurrence. To
support multiple unpacking operations, the CPython compiler was updated
to emit specialized group unpacking bytecodes: *
\textbf{\passthrough{\lstinline!BUILD\_LIST\_UNPACK!}}: Pops multiple
sequences off the stack, converts them to tuples/lists, and concatenates
them into a single list. *
\textbf{\passthrough{\lstinline!BUILD\_TUPLE\_UNPACK!}}: Merges multiple
popped sequences into a tuple. *
\textbf{\passthrough{\lstinline!BUILD\_SET\_UNPACK!}}: Merges sequences
into a set, deduplicating elements. *
\textbf{\passthrough{\lstinline!BUILD\_MAP\_UNPACK!}}: Pops multiple
dictionaries/mappings from the stack and merges them into a single
dictionary.

\emph{Note on Compilation Evolution}: In later versions of CPython
(e.g., Python 3.9+), these unpacking instructions were replaced by
highly specialized, lower-overhead instructions such as
\passthrough{\lstinline!LIST\_EXTEND!},
\passthrough{\lstinline!DICT\_MERGE!}, and
\passthrough{\lstinline!DICT\_UPDATE!} to avoid the temporary list
creation overhead for every unpacked component.

\hypertarget{detailed-bytecode-traces}{%
\subsubsection{2. Detailed Bytecode
Traces}\label{detailed-bytecode-traces}}

\hypertarget{case-a-list-unpacking-with-multiple-iterables}{%
\paragraph{Case A: List Unpacking with Multiple
Iterables}\label{case-a-list-unpacking-with-multiple-iterables}}

Consider compiling this list literal containing mixed values and
unpacking:

\begin{lstlisting}[language=Python]
[1, *[2, 3], 4]
\end{lstlisting}

The Python 3.5 compiler generates the following bytecode:

\begin{lstlisting}
  1           0 LOAD_CONST               0 (1)
              3 BUILD_LIST               1              /* Push list [1] onto stack */
              6 LOAD_CONST               1 (2)
              9 LOAD_CONST               2 (3)
             12 BUILD_LIST               2              /* Push list [2, 3] onto stack */
             15 LOAD_CONST               3 (4)
             18 BUILD_LIST               1              /* Push list [4] onto stack */
             21 BUILD_LIST_UNPACK        3              /* Merge the 3 list elements on the stack */
\end{lstlisting}

\passthrough{\lstinline!BUILD\_LIST\_UNPACK 3!} pops the three
constructed list elements off the stack, performs dynamic sequence
iteration to concatenate their values, and pushes the final unified list
\passthrough{\lstinline![1, 2, 3, 4]!} back onto the evaluation stack.

\hypertarget{case-b-dictionary-unpacking-with-multiple-mappings}{%
\paragraph{Case B: Dictionary Unpacking with Multiple
Mappings}\label{case-b-dictionary-unpacking-with-multiple-mappings}}

Consider this dictionary literal merging two mappings:

\begin{lstlisting}[language=Python]
y = {**d1, **d2, 'key': 42}
\end{lstlisting}

The Python 3.5 compiler generates the following bytecode:

\begin{lstlisting}
  1           0 LOAD_FAST                0 (d1)         /* Push dictionary d1 */
              3 LOAD_FAST                1 (d2)         /* Push dictionary d2 */
              6 LOAD_CONST               2 ('key')
              9 LOAD_CONST               3 (42)
             12 BUILD_MAP                1              /* Push dictionary {'key': 42} */
             15 BUILD_MAP_UNPACK         3              /* Merge the 3 dict elements on the stack */
             18 STORE_FAST               2 (y)
\end{lstlisting}

During execution: * \passthrough{\lstinline!BUILD\_MAP\_UNPACK 3!}
evaluates the three mappings popped from the stack. * It iterates
through their keys, updating the target dictionary. * In Python 3.5, key
collisions do not raise an error for dictionary literals; rather, keys
to the right overwrite keys to the left (e.g.,
\passthrough{\lstinline!d2!} overrides \passthrough{\lstinline!d1!}, and
\passthrough{\lstinline!'key'!} overrides any prior matching key). * In
contrast, if multiple unpacking overlaps occur during keyword argument
calls (e.g.~\passthrough{\lstinline!func(**d1, **d2)!}), duplicate keys
trigger a runtime \passthrough{\lstinline!TypeError!} exception.

\begin{center}\rule{0.5\linewidth}{0.5pt}\end{center}
